\SourceLeaf{380}%
for a time, and there was a considerable variation in the collimation error, the like of which had \textit{never} occurred before. Its cause was first sought in vain in a number of other circumstances. It turned out to be a little knot in the thread tied around the bladder membrane that kept the spirit level sealed. The knot was just a little too thick and touched the inner surface of the sleeve. When it was cut away, so that the glass tube touched only the slightly protruding tips of the screws, the collimation error became constant again and the earlier excellent agreement of all the observations was restored at once. In the Turin observations under discussion, too, we notice considerable variation in the collimation error. We mean the smaller changes, which appear to be accidental, and leave aside the larger changes of several minutes, which no doubt had a good reason known to the astronomer, though it is unmentioned in the observations. Now that there is such an excellent means of determining the collimation error at any moment without reversing the instrument, the cause of similar anomalies can be found all the more easily.

\begingroup\rightskip=0pt plus 1em
The fourth section describes how Mont Cenis was tied into the triangulation network by a separate triangulation and a short baseline measured on the mountain plateau. It also gives the astronomical determination of the azimuth of the line Mont Cenis --- Bellecombe. This determination was made twice. The results for 1821, 1822, obtained with repeating circles by \textsc{Throughton} and \textsc{Reichenbach}, differ by \(8''6\), and the mean of the two was taken, although the later determination seems far more reliable.\par
\endgroup

Similarly, the following two sections contain the astronomical determinations of the azimuths of the lines Mont Colombier --- Mont Granier and Turin new observatory --- Superga. In all three cases the meridian marks of the respective transit telescopes were the basis for these determinations. Finally, the sixth section also gives the operations by which an astronomical determination of azimuth that \textsc{Oriani} had made earlier at the Milan observatory was transferred to the orientation of the side Milan cathedral tower --- Busto in the French triangulation network.

The seventh section then derives the results for the measurement of the arc of longitude from these extensive operations. The measurements were referred
